\documentclass[10pt,a4paper]{article}
\usepackage{oneorder-theme}
\renewcommand{\reportRound}{Round 3}
\newcommand{\chipnew}{\chip{orangetint}{risktx}{New in Round 3}}
\begin{document}
\masthead{03 · The solution · User journey}{User journey}

\htwo{1. Persona and scenario}
\begin{oocard}[accentleft]
\renewcommand{\arraystretch}{1.1}\small
\begin{tabularx}{\linewidth}{@{}>{\bfseries\raggedright\arraybackslash\leavevmode\color{navy}}p{2.2cm}Y@{}}
Who & \textbf{Anna}, sales desk employee at Company A in Denmark, handling B2B project orders all day. \\
Trigger & \textbf{Mette Hoffmann} of \textbf{Hoffmann Byg A/S} emails a material list for the harbour site at \emph{Havnegade 45, 5000 Odense C} with a fixed delivery date. The same journey works for a photo of a handwritten list, a PDF or an Excel file. \\
Pain today & Anna retypes every line, sees the stock of one country only, cannot promise a date she trusts, and keys the customer into the CRM separately. \\
Goal & One tool: a confirmed order with reserved stock, a safe delivery date, and the order on the right customer in Salesforce. \\
Account owner & \textbf{Freja} manages Hoffmann Byg A/S in Salesforce and sees the confirmed project order there without opening OneOrder. \chipnew \\
Year-1 scope & Sweden (no CRM, customer typed as before) and Denmark (Salesforce from month 3). Norway joins in year 2. \chipnew \\
\end{tabularx}
\end{oocard}

\htwo{2. The journey at a glance}
\begin{center}
\begin{tikzpicture}[scale=0.98,transform shape]
\tikzset{jstep/.style={ooc,minimum width=2.45cm,minimum height=2.15cm,inner sep=4pt,align=center}}
\def\g{2.99}
\node[jstep,bluetop]   (a) at (0*\g,0) {{\sgfont\fontsize{14}{14}\selectfont\color{orange}1}\\[1pt]\textbf{Receive}\\Email, photo,\\PDF or Excel\\arrives};
\node[jstep,bluetop]   (b) at (1*\g,0) {{\sgfont\fontsize{14}{14}\selectfont\color{orange}2}\\[1pt]\textbf{AI reads}\\Draft order,\\confidence\\per line};
\node[jstep,accenttop] (c) at (2*\g,0) {{\sgfont\fontsize{14}{14}\selectfont\color{orange}3}\\[1pt]\textbf{Anna reviews}\\Lines, and picks\\the customer\\from Salesforce};
\node[jstep,bluetop]   (d) at (3*\g,0) {{\sgfont\fontsize{14}{14}\selectfont\color{orange}4}\\[1pt]\textbf{Plan}\\Stock in DK\\and SE; shipments\\per source};
\node[jstep,bluetop]   (e) at (4*\g,0) {{\sgfont\fontsize{14}{14}\selectfont\color{orange}5}\\[1pt]\textbf{ETA}\\P50 / P90 date\\per shipment\\vs.\ request};
\node[jstep,accenttop] (f) at (5*\g,0) {{\sgfont\fontsize{14}{14}\selectfont\color{orange}6}\\[1pt]\textbf{Confirm}\\Reserves stock,\\then syncs the\\order to Salesforce};
\foreach \x/\y in {a/b,b/c,c/d,d/e,e/f} \draw[flow] (\x.east) -- (\y.west);
\node[font=\semibold\fontsize{7.5}{9}\selectfont,text=risktx,fill=orangetint,rounded corners=3pt,inner sep=2pt] at ([yshift=-11pt]c.south) {new: CRM customer};
\node[font=\semibold\fontsize{7.5}{9}\selectfont,text=risktx,fill=orangetint,rounded corners=3pt,inner sep=2pt] at ([yshift=-11pt]f.south) {new: CRM sync};
\end{tikzpicture}
\end{center}
\vspace{-4pt}
\figurenote{Orange edge = a decision by a person. Blue edge = done by the system. A person always has the last word before stock is reserved. In year 1 the Norwegian ERP is not connected; the MVP demo still shows all three countries.}

\htwo{3. Main journey step by step}
{\footnotesize
\renewcommand{\arraystretch}{1.12}
\newlength{\jw}\setlength{\jw}{\dimexpr(\linewidth-0.5cm-2.4cm-6\tabcolsep)/2\relax}
\begin{longtable}{@{}>{\bfseries\raggedright\arraybackslash\leavevmode\color{orangetext}}p{0.5cm}>{\bfseries\raggedright\arraybackslash\leavevmode\color{navy}}p{2.4cm}>{\raggedright\arraybackslash}p{\jw}>{\raggedright\arraybackslash}p{\jw}@{}}
\# & Step & What Anna does & What OneOrder does \\ \thead \endfirsthead
\# & Step & What Anna does & What OneOrder does \\ \thead \endhead
1 & Receive the order & Opens the \textbf{Sample inbox} (or uploads a file, or pastes the email text) and imports Mette's message. & Creates a \emph{draft} order and keeps the original file next to the result. \\ \trow
2 & AI extraction & Sees the new order under \emph{To review} on the \textbf{Dashboard}. & Reads customer, phone, delivery site, requested date and every line. Matches each line to the product catalogue. Lines under \textbf{85\%} confidence, or where fuzzy matching and AI disagree, are flagged. \\ \trow
\hl 3 & \hl Review and approve & \hl Compares the original (left) with what the system read (right), fixes flagged lines and ticks them. Picks \textbf{Hoffmann Byg A/S} from the Salesforce suggestions (matched on +45 33 12 45 78). If several accounts share the number, she checks name and city first. If the customer is not in the CRM, she types it. Presses \emph{Approve}. & \hl Blocks \emph{Approve} until every flagged line is checked and a customer with a valid phone is present. Suggests CRM accounts by phone and name and warns when the match is ambiguous; it never links a customer on its own. Editing the name or number later removes the link. Every change is written to the \textbf{audit log}. \\ \trow
4 & Plan fulfilment & Keeps the detected site (Odense) and presses \emph{Plan}. & Reads stock from the national ERPs through adapters (Denmark and Sweden in year 1). Picks sources by rule: nearest branch → national DC → DC in another country. Splits into shipments with weight, vehicle and cross-border flag, states \textbf{why} each source was chosen and lists shortages. \\ \trow
5 & See the delivery date & Reads the \textbf{ETA} per shipment and for the whole order. & Predicts a \textbf{P50} (expected) and a \textbf{P90} (safe to promise) date and compares it with Mette's date: \chipok{on time}, \chiprisk{at risk} or \chiplate{late}. \\ \trow
\hl 6 & \hl Confirm & \hl Presses \emph{Confirm}. A shortage or a late risk needs an explicit tick. & \hl Reserves the planned stock in the ERPs and freezes plan and ETA in one transaction. \textbf{Only then} writes the order to Salesforce, keyed by the OneOrder order id so a retry never creates a duplicate. Shows \emph{Synced to Salesforce}; Freja sees the order on the account. \\
\end{longtable}}

\htwo{4. Moments where the journey branches}
{\footnotesize
\renewcommand{\arraystretch}{1.1}
\begin{tabularx}{\linewidth}{@{}>{\bfseries\raggedright\arraybackslash\leavevmode\color{navy}}p{4.4cm}Y@{}}
Situation & What happens \\ \thead
Low-confidence or unmatched line & \emph{Approve} stays blocked until Anna fixes or confirms the line. \\ \trow
Customer name not detected & The customer field is required; Anna types it or picks a CRM suggestion. \\ \trow
\hl Several CRM accounts share the phone & \hl All candidates are shown with a warning to check carefully; Anna decides by name and city. \\ \trow
\hl Customer not in the CRM, or the lookup fails & \hl Anna types the customer; the order is not linked and shows \emph{Not synced: customer not linked}. \\ \trow
\hl Order for a Swedish site & \hl Sweden has no CRM in year 1: the customer is typed as before and the order shows \emph{Not synced: no CRM for this country}. \\ \trow
\hl Salesforce does not respond at Confirm & \hl Confirm succeeds, stock stays reserved. \emph{Pending sync}, retried after 30\,s and 120\,s; after 3 attempts \emph{Failed} with the reason and \emph{Retry sync}. \\ \trow
Stock short, or taken by another order & The plan lists the shortage; Confirm needs an explicit partial-order tick, or Anna plans again. \\ \trow
Predicted arrival after the requested date & The order is marked \chiplate{late}. Confirm needs \emph{accept late risk}; Anna can call Mette with the P90 date instead. \\ \trow
Sourcing or CRM policy changes & A lead edits \emph{Rules} (sources, cross-border, CRM per country); no redeploy. \\
\end{tabularx}}

\htwo{5. What changed since Round 2}
{\footnotesize
\renewcommand{\arraystretch}{1.1}
\begin{tabularx}{\linewidth}{@{}>{\bfseries\raggedright\arraybackslash\leavevmode\color{navy}}p{3.1cm}>{\raggedright\arraybackslash}p{3.5cm}Y@{}}
Step or element & Round 2 & Round 3 and why \\ \thead
Countries in year 1 & DK, SE, NO & SE and DK; NO in year 2, which funds the CRM work inside the 30\% cut. \\ \trow
Customer at review & Name and phone typed & Denmark: picked from Salesforce suggestions, a person decides. Sweden: typed as before. \\ \trow
Confirm & Reserve stock, freeze ETA & Same, then sync the order to Salesforce so the account owner sees it. \\ \trow
Failure handling & Stock conflicts, late risk & Plus: a CRM outage never blocks a confirmation (queue, retry). \\ \trow
People & Anna & Plus Freja, the account owner, who works in Salesforce. \\ \trow
Legacy CRM & Not in scope & Dynamics 365 is never connected to OneOrder; its data moves once to Salesforce (months 0–3). \\
\end{tabularx}}

\vspace{4pt}
\begin{statement}
Anna commits a P90 date with stock reserved; in Denmark Freja sees the order on the right Salesforce customer, and a CRM outage never stops a confirmation.
\end{statement}
\figurenote{Time saved and fewer ``where is my order?'' calls are the team's hypotheses for the pilot, not measured results.}
\end{document}
